\section{ขอบเขตและข้อจำกัดของการวิจัย}

% กำหนดขอบเขตและข้อจำกัดแยกตามแต่ละขั้นตอนของ Data Science Process

\subsection{ขอบเขตแยกตามขั้นตอน}

\begin{table}[H]
\centering
\begin{tabularx}{\textwidth}{|l|L|}
\hline
\textbf{ขั้นตอน} & \textbf{ขอบเขต} \\
\hline
การตั้งคำถาม & ตอบคำถาม 5 ข้อตามที่ระบุไว้ เน้นความสัมพันธ์ระหว่างพฤติกรรมกับสุขภาวะทางจิต ไม่ครอบคลุมการหาสาเหตุ \\
\hline
การเก็บข้อมูล & นักศึกษาระดับปริญญาตรี มหาวิทยาลัยขอนแก่น จำนวน 150--200 คน เก็บด้วยแบบสอบถามออนไลน์เป็นเวลา 14 วัน \\
\hline
การสำรวจข้อมูล & สถิติเชิงพรรณนา การวิเคราะห์ตัวแปรเดี่ยวและหลายตัวแปร พร้อมการตรวจสอบคุณภาพและที่มาของข้อมูล \\
\hline
การวิเคราะห์เชิงสถิติ & การวิเคราะห์สหสัมพันธ์ การเปรียบเทียบค่าเฉลี่ยระหว่างกลุ่ม และการถดถอยพหุคูณ ไม่ครอบคลุมการสร้างแบบจำลองการเรียนรู้ของเครื่อง \\
\hline
ผลลัพธ์ & รายงานผลการวิเคราะห์และการอภิปราย ไม่มีการพัฒนาระบบหรือเครื่องมือสำหรับผู้ใช้ \\
\hline
\end{tabularx}
\end{table}

\subsection{ข้อจำกัด}

\begin{subsubitems}
    \item ผลที่ได้เป็นความสัมพันธ์ ไม่ใช่ความเป็นเหตุเป็นผล เนื่องจากเก็บข้อมูลเพียงครั้งเดียว จึงไม่สามารถระบุได้ว่าการใช้โซเชียลมีเดียทำให้สุขภาวะแย่ลง หรือผู้ที่สุขภาวะไม่ดีอยู่แล้วใช้โซเชียลมีเดียมากกว่า
    \item กลุ่มตัวอย่างเลือกแบบตามสะดวก ไม่ใช่การสุ่มตัวอย่าง จึงไม่สามารถอ้างอิงไปยังนักศึกษาทั้งมหาวิทยาลัยหรือทั้งประเทศได้
    \item ข้อมูลบางส่วนเป็นการรายงานตนเอง คะแนนในสเกล 1--10 ของแต่ละคนอาจไม่ได้หมายถึงระดับความรู้สึกที่เท่ากัน แม้จะลดปัญหานี้ในข้อชั่วโมงการใช้งานด้วยการให้อ่านค่าจากเครื่องแล้วก็ตาม
    \item ขนาดกลุ่มตัวอย่างจำกัด หากความสัมพันธ์จริงมีขนาดเล็ก การศึกษานี้อาจตรวจไม่พบ
    \item ผลการศึกษาใช้ได้ในระดับกลุ่ม ไม่ใช่ระดับบุคคล ไม่สามารถนำไปประเมินหรือวินิจฉัยสุขภาพจิตของบุคคลใดบุคคลหนึ่งได้ ผู้ที่มีข้อกังวลควรปรึกษาผู้เชี่ยวชาญ หรือติดต่อสายด่วนสุขภาพจิต 1323
\end{subsubitems}

\subsection{ประเด็นด้านจริยธรรม}

\begin{subsubitems}
    \item ผู้ตอบต้องอ่านคำชี้แจงและกดยืนยันความยินยอมก่อนเข้าสู่แบบสอบถาม
    \item ไม่เก็บชื่อ รหัสนักศึกษา อีเมล หรือข้อมูลใดที่ระบุตัวบุคคลได้
    \item ผู้ตอบสามารถหยุดตอบได้ตลอดเวลาโดยไม่มีผลกระทบใด ๆ
    \item แสดงสายด่วนสุขภาพจิต 1323 ทั้งต้นและท้ายแบบสอบถาม
    \item เลือกใช้แบบวัด WHO-5 ซึ่งไม่มีข้อคำถามเกี่ยวกับการทำร้ายตนเอง
    \item ข้อมูลจะรายงานในภาพรวม และใช้เพื่อการศึกษาในรายวิชาเท่านั้น
\end{subsubitems}
